\section*{Khai báo sử dụng AI và đóng góp}
\addcontentsline{toc}{section}{Khai báo sử dụng AI và đóng góp}

AI có bốn vai trò tách riêng trong đề tài. \textit{Trợ lý nghiên cứu} (Claude Code, Claude Opus 5.5) tìm tài
liệu, soạn nháp đề bài và hai prompt quy trình, viết mã harness, grader, phân tích và script sinh số liệu, viết
nháp báo cáo và slide; An quyết định thiết kế, phạm vi, tiêu chí và mọi thay đổi so với preregistration, và duyệt
đề bài. Mọi con số trong báo cáo sinh bằng máy từ file kết quả. \textit{Người duyệt độc lập} (Fable 5.2, ngữ
cảnh riêng) phản biện preregistration, đề bài và mã phân tích trước khi có kết quả, chấm checklist của E2 khi
không biết cách làm, và gán nhãn thất bại như người thứ ba chỉ để đối chiếu. \textit{Đối tượng được đo} là các
cohort ở mục~\ref{sec:conditions}; trợ lý nghiên cứu cùng mô hình với \texttt{claude-opus} và đã soạn đề cho nhóm
này (mục~\ref{sec:threats}). Cuối cùng, AI là công cụ đã làm ra ViPharmind. Nhãn dùng trong phân tích, chấm tay
mẫu E2 và đánh giá chéo do người làm, không dùng AI.

\textit{Đóng góp.} Doãn Văn An (nhóm trưởng) điều phối và làm phần lớn nghiên cứu: thiết kế, bộ nhiệm vụ,
harness, chạy thí nghiệm, phân tích, báo cáo, slide. Lê Xuân Lợi và Trần Ngọc Bảo cùng góp ý thiết kế và làm
độc lập phần đo bằng người (gán nhãn mù, kiểm tra lượt đạt, chấm tay mẫu E2); Lợi kiểm trích dẫn và phản biện
mục 1--3, Bảo tính lại độc lập tỉ lệ đạt, token và phản biện mục 4--5.
